\section{Introduction}\label{sec:intro}

Phishing is still the most common initial-access vector. Its production cost has collapsed since instruction-tuned language models began to write fluent, personalised lures from a few lines of reconnaissance \cite{heiding2024,hazell2023}; in controlled studies, fully automated LLM spear-phishing matched or exceeded the click-through rates of expert-written messages \cite{heiding2026,bethany2025}, and generator--critic loops now rewrite a lure until machine defenders stop flagging it \cite{qi2025}. Two consequences matter for a detector deployed at an e-mail gateway. The attacker can change the surface form of a campaign at negligible cost, by switching the generating model or by paraphrasing a known lure; that change reaches the gateway as an unannounced shift of the message distribution, observed at full throughput, long before verified labels for the new campaign exist.

Transformer detectors are evaluated almost exclusively on fixed partitions. Fine-tuned DistilBERT and RoBERTa encoders report F1 scores close to 0.99 in that setting \cite{jamal2024,uddin2026,ibrahim2026,safran2025,cetlenbik2024roberta}, yet the last two years of evidence show that such figures do not transfer: LLM rephrasing of known phishing lowers the accuracy of commercial filters and ML detectors \cite{afane2024,thapa2025}, and changing the generating model between training and deployment costs 28 F1 points while the AUC stays above 0.96 \cite{gutierrez2026}. The ranking survives that change; the operating threshold does not. Template artefacts in generated corpora can also turn in-distribution scores into a ceiling effect rather than evidence of learned semantics \cite{cetlenbik2026}.

Prior work leaves the streaming case open. The studies that model generative-AI shift explicitly do so offline, on predefined partitions \cite{arharbi2026,eskandarian2026}; the drift literature offers label-free detectors for deep representations \cite{rabanser2019,greco2025,wan2024,feldhans2024} and mature adaptation loops for malware and intrusion telemetry \cite{pan2026,han2026,yang2025,mcfadden2026}, but none of them has been specified for e-mail streams in which labels arrive late and cost analyst time while the shift itself may be adversarial. This paper asks whether a transformer-based phishing detector can notice, without labels, that the message distribution has moved, and then recover its accuracy from a small purchased label budget while paraphrase campaigns drive the move. We answer with SLIDE, a sliced-Wasserstein drift-adaptive defence, and with a prequential protocol built entirely from public corpora.

The contributions are four.
\begin{enumerate}
\item \textbf{A streaming threat model for LLM-era phishing} (Section~\ref{sec:formulation}). It couples delayed, budgeted labels with virtual drift measured by $SW_1$ and KS statistics in representation space, and with evasion families whose budgets are measured rather than assumed. A five-segment public benchmark stream moves from human corpora of 1999--2008 to two independent LLM corpora before ending in a paraphrase campaign; every source is audited for template leakage.
\item \textbf{A label-free sliced two-gate drift test} (Section~\ref{sec:test}). Holm-corrected KS tests on random projections and on the detector score are combined with a sliced-Wasserstein effect-size gate calibrated on the reference window. Proposition~\ref{prop:fwer} proves family-wise false-alarm control; Proposition~\ref{prop:sw} states what an alarm certifies about the underlying shift.
\item \textbf{A label-on-alarm adaptation policy with a composite robust objective} (Sections~\ref{sec:loss} and~\ref{sec:policy}). The objective has four terms: class-weighted cross-entropy on the purchased window, class-balanced reservoir replay, an $\ell_2$ adversarial loss in representation space and a proximal penalty that limits each update. A Lipschitz margin certificate (Proposition~\ref{prop:cert}) and a time and memory analysis (Proposition~\ref{prop:cost}) accompany the objective.
\item \textbf{A statistically controlled empirical study} (Section~\ref{sec:results}). SLIDE is compared with nine baselines, from static detectors to fully supervised retraining, over ten stream realisations under Friedman/Nemenyi and Holm-corrected Wilcoxon tests. The study adds three equal-budget timing controls, a nine-variant ablation, a sensitivity analysis, a six-variant evasion stress test and component-level latency on commodity CPU hardware; all code is released.
\end{enumerate}
The remainder follows the order of these contributions, with the results in Section~\ref{sec:results} and their limits in Section~\ref{sec:discussion}.
